\chapter{ניהול סיכונים}\label{ch:risks}

סיכון הוא הסתברות לאירוע כפול ההשפעה שלו. RoofRoot מנהלת שמונה סיכונים מרכזיים ברישום אחד, ומדרגת כל אחד מהם בסולם של אחד עד חמש לפי סבירות והשפעה. הרישום מופיע בטבלה~\ref{tab:risks} והמיקום של כל סיכון על מפת החום מוצג באיור~\ref{fig:risk-map}. שלושה סיכונים נמצאים באזור הגבוה ביותר ומקבלים תשומת לב מיוחדת בהמשך הפרק.

\section{מפת הסיכונים}

מפת החום ממקמת כל סיכון לפי סבירות (ציר אופקי) והשפעה (ציר אנכי). הסיכונים הקרובים לפינה העליונה דורשים מיתון פעיל ובעלים בהנהלה. אלה שבפינה התחתונה נבדקים אחת לרבעון.

\begin{figure}[htbp]
\centering
\includegraphics[width=0.85\linewidth]{fig_risk_map}
\caption{מפת חום של הסיכונים. הסיכונים \en{R1} ו-\en{R8} בעלי ההשפעה הגבוהה ביותר, ו-\en{R3} הסבירות הגבוהה ביותר.}
\label{fig:risk-map}
\end{figure}

\begin{table}[htbp]
\centering\small
% AUTO-GENERATED by figures/python/make_figures.py — do not edit by hand.
\rrstripes
\begin{tabularx}{\linewidth}{>{\raggedright}p{6mm}>{\raggedright\arraybackslash}p{3.4cm}cc>{\raggedright\arraybackslash\small}X}
\rrtablehead \rrth{\en{ID}} & \rrth{סיכון} & \rrth{ס.} & \rrth{ה.} & \rrth{מיתון} \\
R1 & מגבלות עומס של הגג & 3 & 5 & סקר מהנדס לפני כל חוזה; מערכת \en{NFT} קלה (פחות מ-60 ק"ג למ"ר) \\
R2 & החלפת ועד בית & 3 & 3 & חוזה ל-5 שנים וחלוקת הכנסות שמתגמלת את הדיירים \\
R3 & לחץ על מחירי התוצרת & 4 & 3 & תמהיל עשיר בעשבי תיבול; חוזי \en{B2B} ישירים עם מסעדות \\
R4 & התפרצות מזיקים או מחלה & 2 & 4 & מערכת סגורה ומודל ראייה של \en{RootOS} להתרעה מוקדמת \\
R5 & גלי חום והפסקות חשמל & 3 & 4 & רשתות הצללה וסוללת גיבוי למשאבות \\
R6 & עיכוב רגולטורי בהיתרים & 2 & 3 & פיילוט בשיתוף יחידת האקלים העירונית \\
R7 & תלות באנשי מפתח & 2 & 4 & נהלים כתובים; \en{RootOS} משמר את הידע של המגדלים \\
R8 & פער מימון לפני סבב \en{A} & 3 & 5 & סבב \en{A} במנות לפי אבני דרך; הכנסות \en{SaaS} מ-2029 \\
\end{tabularx}

\caption{רישום הסיכונים. ס.\ = סבירות, ה.\ = השפעה, בסולם 1--5 (1 נמוך, 5 גבוה). עמודת המיתון מתארת את הפעולה המרכזית.}
\label{tab:risks}
\end{table}

\section{שלושת הסיכונים המרכזיים}

\subsection{\texorpdfstring{\en{R1}}{R1}: מגבלות מבניות}

גג שאינו עומד בעומס אינו יכול לשמש להתקנה, ולכן זה הסיכון בעל ההשפעה הגבוהה ביותר, יחד עם מימון. הסבירות בינונית: בניינים ישנים רבים לא תוכננו לעומס נוסף. מערכת \en{NFT} קלה היא פתרון חלקי, כי היא מפחיתה את העומס, אך אינה מבטלת את הצורך בבדיקה.

\begin{riskbox}[R1 - מגבלות מבניות]
\textbf{מיתון.} סקר מהנדס קונסטרוקציה לפני כל חוזה, ומערכת קלה במשקל נמוך מ-60 ק"ג למ"ר. גג שנכשל בסקר אינו נכנס לצי, והעלות של הסקר נכללת בהשקעה לגג, \RRcapexRoof{} אלף \shekel{}. \textbf{מחיר הכישלון.} תקלה אחת בבניין תפגע במוניטין בפני כל ועדי הבתים, ולכן אין פשרות בבדיקה.
\end{riskbox}

\subsection{\texorpdfstring{\en{R8}}{R8}: פער במימון לפני \texorpdfstring{\en{Series A}}{Series A}}

לפי עקומת המזומנים (איור~\ref{fig:cash-curve}) היתרה מגיעה לשפל של \RRminCash{} מיליון \shekel{} ב-\RRminCashYear{}. עיכוב בגיוס \en{Series A} של \RRseriesA{} מיליון \shekel{} או עלייה בעלויות עלולים להביא את החברה לחוסר נזילות לפני שהגיעה לכלכלה יציבה.

\begin{riskbox}[R8 - פער במימון]
\textbf{מיתון.} הסבב מועבר בתשלומים לפי אבני דרך: מספר גגות פעילים, תפוקה למ"ר והכנסה חודשית. הכנסות \en{SaaS} מ-2029 מקטינות את התלות בגיוס. בנוסף, ההנהלה יכולה להאט את פריסת הגגות, כי הוצאות ההקמה הן ההוצאה הגמישה ביותר. \textbf{הסימן המקדים.} אם יתרת המזומן יורדת מתחת לרצפה שקבע הדירקטוריון, הפריסה נעצרת עד להבהרה.
\end{riskbox}

\subsection{\texorpdfstring{\en{R3}}{R3}: לחץ על מחירי התוצרת}

זה הסיכון הסביר ביותר, כי מחירי ירקות עלים משתנים עונתית ומושפעים מיבוא. כפי שראינו בפרק~\ref{ch:financials}, מחיר ותפוקה הם הגורם הדומיננטי ב-\en{NPV}, ולכן ירידת מחיר מורגשת בכל ההנחות.

\begin{riskbox}[R3 - לחץ על מחירים]
\textbf{מיתון.} תמהיל עשיר בעשבי תיבול, שמחירם גבוה יותר ויציב יותר, וחוזים ישירים עם מסעדות שמקבעים מחיר לתקופה. מנוי \enquote{סל ירוק} לתושבים מוסיף ערוץ שבו המחיר פחות תחרותי. \textbf{בקרה.} המחיר הממוצע לק"ג נבדק חודשית מול ההנחה של \RRprice~\shekel.
\end{riskbox}

\section{סיכונים נוספים}

חמשת הסיכונים הנותרים נמצאים בטווח בינוני ונוהלו במיתון סטנדרטי. עזיבת ועד בית (\en{R2}) מטופלת בחוזה ל-\RRcontractYears{} שנים שבו התושבים חולקים בהכנסה. מחלה במערכת (\en{R4}) מזוהה מוקדם על ידי מודל ראיית העלה של \en{RootOS}. גל חום או הפסקת חשמל (\en{R5}) מטופלים ברשתות הצללה ובסוללת גיבוי למשאבות. עיכוב ברישוי (\en{R6}) מצטמצם בפיילוט עם יחידת האקלים העירונית. תלות באנשי מפתח (\en{R7}) מצטמצמת בנהלים כתובים, וידע המגדל נשמר גם במערכת.

\section{תהליך ממשל סיכונים}

הרישום הוא כלי חי. התהליך מוגדר כך:

\begin{enumerate}
  \item \textbf{בעלות.} לכל סיכון בעלים בהנהלה: מנכ"לית לסיכוני מימון ומחיר, מנהל טכנולוגיות לסיכוני תוכנה ומערכת, סמנכ"לית תפעול לסיכונים מבניים ותפעוליים.
  \item \textbf{עדכון רבעוני.} כל רבעון מעדכנים סבירות והשפעה ובוחנים אם המיתון עובד. התוצאות מוצגות בישיבת הדירקטוריון (פרק~\ref{ch:team}).
  \item \textbf{אירועים חדשים.} סיכון חדש נכנס לרישום מיד כשזוהה, ולא ממתינים לרבעון הבא.
  \item \textbf{הסלמה.} סיכון שעולה לדירוג גבוה, ציון ארבע ומעלה בשני הצירים, מועבר לדירקטוריון עם תוכנית פעולה בתוך שבועיים.
\end{enumerate}

\begin{insightbox}[המשמעות למשקיע]
שלושת הסיכונים העיקריים אינם סודיים אלא מנוהלים: לכל אחד מיתון, בעלים ומדד שאפשר לעקוב אחריו. הסיכון המבני נבדק לפני החוזה, פער המימון נבדק אחת לרבעון והמחיר נבדק חודשית.
\end{insightbox}
